\section*{अध्याय \ref{ch:b2:plant-life-cycles} --- स्थलीय पौधों के जीवन-चक्र और जनन}

\begin{solution}{exo:b2:plant-life-cycles:1}
\omterm{def:b2:plant-life-cycles:cycles}{युग्मकोद्भिद}: \omterm{def:b2:meiosis-heredity:meiosis}{अगुणित} \omterm{def:b2:unicellular-diversity:unicellular}{बहुकोशिकीय} पीढ़ी, जो किसी बीजाणु से
समसूत्री विभाजन द्वारा उगती है और समसूत्री विभाजन से युग्मक बनाती है।
\omterm{def:b2:plant-life-cycles:cycles}{बीजाणुद्भिद}: \omterm{def:b2:meiosis-heredity:meiosis}{द्विगुणित} पीढ़ी, जो किसी युग्मनज से समसूत्री विभाजन द्वारा
उगती है और \omterm{def:b2:meiosis-heredity:meiosis}{अर्धसूत्री विभाजन} से बीजाणु बनाती है।
बीजाणु: \omterm{def:b2:meiosis-heredity:meiosis}{अर्धसूत्री विभाजन} से बनी ऐसी \omterm{def:b2:meiosis-heredity:meiosis}{अगुणित} कोशिका जो बिना संलयन के
उगती है। युग्मक: समसूत्री विभाजन से (पौधों में) बनी ऐसी \omterm{def:b2:meiosis-heredity:meiosis}{अगुणित} कोशिका
जिसे किसी दूसरी से संलयित होना ही पड़ता है।
\end{solution}

\begin{solution}{exo:b2:plant-life-cycles:2}
पर्ण-कोशिकाएँ $2n = 1260$ हैं: बीजाणु 630, \omterm{def:b2:plant-life-cycles:fern}{प्रोथैलस} कोशिका 630,
शुक्राणु 630, अंडाणु 630, युग्मनज 1260।
\end{solution}

\begin{solution}{exo:b2:plant-life-cycles:3}
हरा गद्दा: \omterm{def:b2:plant-life-cycles:cycles}{युग्मकोद्भिद}। डंठल और संपुट: \omterm{def:b2:plant-life-cycles:cycles}{बीजाणुद्भिद}। \omterm{def:b2:plant-life-cycles:fern}{फ़र्न}
पत्रक: \omterm{def:b2:plant-life-cycles:cycles}{बीजाणुद्भिद}। \omterm{def:b2:plant-life-cycles:fern}{प्रोथैलस}: \omterm{def:b2:plant-life-cycles:cycles}{युग्मकोद्भिद}। \omterm{def:b2:plant-life-cycles:heterospory}{पराग} कण: नर
\omterm{def:b2:plant-life-cycles:cycles}{युग्मकोद्भिद}। बीज का भोजन-भंडार (किसी चीड़ में): मादा \omterm{def:b2:plant-life-cycles:cycles}{युग्मकोद्भिद}।
बीज का भ्रूण: अगला \omterm{def:b2:plant-life-cycles:cycles}{बीजाणुद्भिद}।
\end{solution}

\begin{solution}{exo:b2:plant-life-cycles:4}
अगुणितिक (\emph{Chlamydomonas}): युग्मनज में तुरंत \omterm{def:b2:meiosis-heredity:meiosis}{अर्धसूत्री विभाजन}।
द्विगुणितिक (प्राणी, \emph{Fucus}): \omterm{def:b2:meiosis-heredity:meiosis}{अर्धसूत्री विभाजन} युग्मक बनाता है।
अगुणित-द्विगुणितिक (\omterm{def:b2:plant-life-cycles:moss}{मॉस}, \omterm{def:b2:plant-life-cycles:fern}{फ़र्न}, बीज-पौधे, अनेक शैवाल): यानी
\omterm{def:b2:meiosis-heredity:meiosis}{अर्धसूत्री विभाजन}, जो
\omterm{def:b2:plant-life-cycles:cycles}{बीजाणुद्भिद} की बीजाणुधानियों में बीजाणु बनाता है।
\end{solution}

\begin{solution}{exo:b2:plant-life-cycles:5}
$0.03/10^{-4} = \qty{300}{s}$, यानी पाँच मिनट। \qty{20}{min} में
\qty{12}{cm}। निषेचन केवल कुछ सेंटीमीटर दूर के पौधों के बीच ही चलता है,
इसलिए \omterm{def:b2:plant-life-cycles:moss}{मॉस} सघन गद्दों में उगते हैं जिनमें दोनों लिंग (या दोनों अंग)
पहुँच के भीतर रहते हैं।
\end{solution}

\begin{solution}{exo:b2:plant-life-cycles:6}
$v_s = 2\times(1.5\times 10^{-5})^2\times 1099\times 9.81/(9\times
1.8\times 10^{-5}) = \qty{0.030}{m/s}$, यानी \qty{3}{cm/s}; दूरी $uh/v_s
= 2\times 0.5/0.03 = \qty{33}{m}$। बीज: $2\times 20/0.8 =
\qty{50}{m}$ --- यानी पंख वाली कोई भारी इकाई ऊँचाई से झड़कर लगभग उतनी
ही दूर जाती है।
\end{solution}

\begin{solution}{exo:b2:plant-life-cycles:7}
प्रति पत्रक $200\times 40\times 64 = \num{512000}$ बीजाणु; $5.1$
\omterm{def:b2:plant-life-cycles:fern}{प्रोथैलस}; प्रति पत्रक $0.1$ नया \omterm{def:b2:plant-life-cycles:cycles}{बीजाणुद्भिद}।
\end{solution}

\begin{solution}{exo:b2:plant-life-cycles:8}
बीज एक रोका हुआ \omterm{def:b2:plant-life-cycles:heterospory}{बीजांड} है: मादा \omterm{def:b2:plant-life-cycles:cycles}{युग्मकोद्भिद} को जनक पर ही
रखना, घेरना और पालना पड़ता है, जबकि नर \omterm{def:b2:plant-life-cycles:cycles}{युग्मकोद्भिद} को
उस तक जाना पड़ता है। इसके लिए दो नियतियों वाले दो प्रकार के बीजाणु
चाहिए। किसी समबीजाणुक पौधे का एकमात्र बीजाणु झड़ना ही चाहिए ताकि वह कोई
स्वतंत्र उभयलिंगी \omterm{def:b2:plant-life-cycles:cycles}{युग्मकोद्भिद} बन सके; उसे रोक लेने से नर कार्य भी रुक
जाता, हर पौधे पर स्व-निषेचन थोप दिया जाता, और यात्रा करने को कुछ
बचता ही नहीं।
\end{solution}

\begin{solution}{exo:b2:plant-life-cycles:9}
बीजाणु: $\tfrac43\pi(1.5\times 10^{-5})^3\times 1100 =
\qty{1.6e-11}{kg}$, यानी \qty{16}{ng}; ऊर्जा $0.5\times 1.6\times
10^{-8}\times 2\times 10^{4} = \qty{1.6e-4}{J}$। बीज: \qty{5}{mg},
$0.5\times 5\times 10^{-3}\times 2\times 10^{4} = \qty{50}{J}$ ---
यानी $3\times 10^{5}$ गुना अधिक। बीजाणु प्रकाश से पहले एक छोटी
प्रकाशसंश्लेषी पपड़ी भर उगा सकता है और कुछ नहीं; बीज अँधेरे में जड़ और
प्ररोह उगा सकता है, एक ऋतु प्रतीक्षा कर सकता है, और खा लिए जाने या सूख
जाने से बच सकता है।
\end{solution}

\begin{solution}{exo:b2:plant-life-cycles:10}
प्रति घन मीटर $10^{11}/(10^{6}\times 20) = \num{5000}$ कण; अभिवाह
$5000\times 10^{-7}\times 2 = \qty{1e-3}{}$ कण प्रति सेकंड; यानी प्रति
कण लगभग \qty{1000}{s}, एक चौथाई घंटा। शंकु घंटों से दिनों तक खुला और
ग्रहणशील रहना चाहिए, और ठीक तभी जब नर शंकु झाड़ रहे हों --- परागण का
समय सप्ताह भर की परिशुद्धता से बैठाया जाता है।
\end{solution}

\begin{solution}{exo:b2:plant-life-cycles:11}
द्विगुणिता \omterm{def:b2:meiosis-heredity:vocabulary}{अप्रभावी} हानिकारक \omterm{def:b2:genome-diversification:mutation}{उत्परिवर्तन} ढँक देती है; संवहनी ऊतक वाली
कोई \omterm{def:b2:meiosis-heredity:meiosis}{द्विगुणित} काया बड़ी और दीर्घजीवी हो सकती है, और ऊँचाई प्रकाश-ग्रहण
तथा \omterm{thm:b2:plant-life-cycles:dispersal}{बीजाणु प्रकीर्णन} दोनों में सहायक है; किसी \omterm{def:b2:plant-life-cycles:cycles}{बीजाणुद्भिद} पर ऊँचे बने
बीजाणु वायु में प्रकीर्णित होते हैं, जबकि युग्मकों को तैरना पड़ता है।
\omterm{def:b2:meiosis-heredity:meiosis}{अगुणित} पीढ़ी इसलिए बनी रही कि \omterm{def:b2:meiosis-heredity:meiosis}{अर्धसूत्री विभाजन} और निषेचन को कोई
\omterm{def:b2:meiosis-heredity:meiosis}{अगुणित} प्रावस्था चाहिए ही, और \omterm{def:b2:plant-life-cycles:heterospory}{पराग} कण तथा भ्रूणकोष वे न्यूनतम काएँ हैं
जो युग्मकों को एक-दूसरे तक पहुँचाती हैं और भ्रूण को पालती हैं।
\end{solution}

\begin{solution}{exo:b2:plant-life-cycles:12}
बीजाणुओं का संवर्धन करके और उन्हें सूक्ष्मदर्शी के नीचे देखकर वे यह
स्थापित कर सके कि कोई बीजाणु उगकर लिंग-अंगों वाली एक छोटी काया बनता है,
कि भ्रूण \omterm{def:b2:plant-life-cycles:moss}{स्त्रीधानी} के भीतर निषेचित अंडाणु से उठता है,
कि बीजाणुधारी पौधा उसी से उगता है, और
कि किसी शंकुधारी के \omterm{def:b2:plant-life-cycles:heterospory}{बीजांड} में वही अंग घटे हुए रूप में हैं: यानी
कायाओं और अंगों के अनुक्रम के रूप में वह एकांतरण। वे यह नहीं जान सकते थे
कि कोशिका के स्तर पर दोनों पीढ़ियों में भेद क्या है।
गुणसूत्र-गणनाओं ने दिखाया कि दोनों कायाओं की गुणसूत्र-संख्या में दो का
गुणक अंतर है, कि \omterm{def:b2:meiosis-heredity:meiosis}{अर्धसूत्री विभाजन} बीजाणु-निर्माण पर बैठता है और
निषेचन अंडाणु पर, और यों समझाया कि वह एकांतरण अनिवार्य क्यों है।
\end{solution}

\begin{solution}{pb:b2:plant-life-cycles:1}
\textbf{1.} $10\times 200\times 40\times 64 = 5.1\times 10^{6}$
बीजाणु।
\textbf{2.} प्रति बीजाणु $\tfrac43\pi(1.5\times 10^{-5})^3\times 1100 =
\qty{1.6e-11}{kg}$; फ़सल $8\times 10^{-5}$ kg, यानी \qty{80}{mg}।
\textbf{3.} $v_s = 2\times 2.25\times 10^{-10}\times 1099\times
9.81/(1.62\times 10^{-4}) = \qty{0.030}{m/s}$।
\textbf{4.} $x = 1.5\times 0.5/0.030 = \qty{25}{m}$;
$\pi\times 25^2 = \qty{2000}{m^2}$ की चक्रिका; यानी प्रति वर्ग मीटर
लगभग \num{2600} बीजाणु।
\textbf{5.} कुछ सेंटीमीटर प्रति सेकंड की ऊर्ध्वधाराएँ $v_s$ से बड़ी
हैं, इसलिए विक्षोभ में फँसा कोई बीजाणु घंटों हवा में रहता है और सैकड़ों
किलोमीटर जाता है; और एक ही बीजाणु कोई आबादी बसा देता है यदि उसका
\omterm{def:b2:plant-life-cycles:fern}{प्रोथैलस} स्व-निषेचन कर सके (वह दोनों अंग ढोता है), और इसीलिए
दूरदराज़ के द्वीपों की \omterm{def:b2:plant-life-cycles:fern}{फ़र्न} वनस्पति समृद्ध है।
\textbf{6.} बीजाणु: $0.5\times 1.6\times 10^{-8}\times 2\times 10^{4} =
\qty{1.6e-4}{J}$; फ़सल: \qty{820}{J}; एक चीड़-बीज: \qty{50}{J} ---
यानी पूरी बीजाणु-फ़सल सोलह बीजों के बराबर है।
\textbf{7.} $5.1\times 10^{6}/2\times 10^{4} = 256$ \omterm{def:b2:plant-life-cycles:fern}{प्रोथैलस}।
\textbf{8.} $1 - 0.8^{6} = 1 - 0.26 = 0.74$।
\textbf{9.} प्रति ऋतु $256\times 0.74\times 0.5\times 0.02 = 1.9$ वयस्क
\omterm{def:b2:plant-life-cycles:fern}{फ़र्न}।
\textbf{10.} $30\times 1.9 \approx 57$: यानी किसी स्थिर आबादी के एक
प्रतिस्थापन से कहीं अधिक, इसलिए उत्तरजीविता के ये आँकड़े आशावादी हैं ---
किसी भरे आवास में अधिकांश नए \omterm{def:b2:plant-life-cycles:cycles}{बीजाणुद्भिद} भीड़ और छाया से मर
जाते हैं।
\textbf{11.} $10^{-3}/10^{-4} = \qty{10}{s}$। अनेक \omterm{def:b2:plant-life-cycles:fern}{फ़र्न}
पुंधानियाँ और स्त्रीधानियाँ भिन्न समयों पर पकाते हैं, या ऐसा हॉर्मोन
(ऐन्थेरिडिओजन) छोड़ते हैं जो पड़ोस के नए \omterm{def:b2:plant-life-cycles:fern}{प्रोथैलसों} को नर बना देता है,
ताकि शुक्राणु किसी दूसरे व्यक्ति से आएँ।
\textbf{12.} वह एक कोशिका मोटा, जड़-रहित, असुरक्षित है, उसे निषेचन के
लिए द्रव जल चाहिए और वह कुछ ही सप्ताह जीता है: चक्र की लगभग सारी
मृत्यु-दर उसी पर पड़ती है ($2\times 10^{4}$ में एक बीजाणु उसमें बदलता
है; और उनमें से एक चौथाई कोई गीली रात कभी नहीं देखते)।
\textbf{13.} $v_s = 2\times 6.25\times 10^{-10}\times 599\times
9.81/(1.62\times 10^{-4}) = \qty{0.045}{m/s}$।
\textbf{14.} $3\times 15/0.045 = \qty{1000}{m}$।
\textbf{15.} प्रति घन मीटर $10^{11}/(2\times 10^{7}) = \num{5000}$।
\textbf{16.} $5000\times 10^{-7}\times 2 = \qty{1e-3}{s^{-1}}$: यानी
प्रति घंटा 3.6; पहला कण लगभग \qty{1000}{s}, यानी एक चौथाई घंटे बाद आता
है।
\textbf{17.} पाँच किलोमीटर अनुवात वह बादल अगल-बगल और ऊपर फैल चुका है और
अधिकांश कण बैठ चुके हैं (प्रति \qty{15}{m} ऊँचाई पर \qty{1}{km}
परास), इसलिए सांद्रता परिमाण की कई कोटि कम है और किसी \omterm{def:b2:plant-life-cycles:heterospory}{बीजांड} को किसी कण
के लिए दिनों प्रतीक्षा करनी पड़ सकती है; और वृक्ष का अपना \omterm{def:b2:plant-life-cycles:heterospory}{पराग}
अधिकतर स्वपरागित बीज देता है, जो नष्ट हो जाते हैं। वायु-परागण को कोई
सघन बादल चाहिए, इसलिए झुंड।
\textbf{18.} $10^{14}$ कण, $2\times 10^{6}$ \omterm{def:b2:plant-life-cycles:heterospory}{बीजांडों} के लिए: प्रति \omterm{def:b2:plant-life-cycles:heterospory}{बीजांड} $5\times
10^{7}$ कण।
\textbf{19.} $2\times 10^{6}$ \omterm{def:b2:plant-life-cycles:heterospory}{बीजांड} $\times 0.6\times 0.8 = 9.6\times
10^{5}$ बीज; \qty{4.8}{kg}; \qty{48}{MJ}।
\textbf{20.} $3\times 20/0.8 = \qty{75}{m}$: \omterm{def:b2:plant-life-cycles:heterospory}{पराग} एक
किलोमीटर जाता है, बीज कुछ ही वृक्ष-ऊँचाइयाँ; जीन \omterm{def:b2:plant-life-cycles:heterospory}{पराग} से
यात्रा करते हैं, पौधा बीज से।
\textbf{21.} $9.6\times 10^{5}\times 0.05\times 0.01 = 480$ वयस्क
संतति, \omterm{def:b2:plant-life-cycles:fern}{फ़र्न} की 57 के मुक़ाबले।
\textbf{22.} चीड़: प्रति वयस्क संतान $4.8\times 10^{7}/480 = \qty{100}{kJ}$;
\omterm{def:b2:plant-life-cycles:fern}{फ़र्न}: प्रति वयस्क संतान $820/1.9 = \qty{430}{J}$, यानी दो
सौ गुना कम।
\textbf{23.} दोनों \qty{20}{kJ/g} पर आधे शुष्क पदार्थ के हैं, इसलिए
मान लीजिए प्रति ग्राम शुष्क \qty{0.5}{mW} की आधारी दर पर दोनों $10^{4}\,
\text{J/g}/(5\times 10^{-4}\,\text{W/g}) = 2\times 10^{7}$ सेकंड, यानी
लगभग आठ महीने चलते हैं: प्रसुप्ति की अवधि वह नहीं है जो बीज की ऊर्जा
ख़रीदती है। वह जो ख़रीदती है वह अंकुरण के बाद होता है --- भंडार से
अँधेरे में गढ़ी गई एक जड़ और एक प्ररोह, अंकुर के स्वयं खाने लगने से
सप्ताहों पहले --- और आवरण के साथ, इस बीच सुरक्षा; जबकि बीजाणु की
\qty{1.6e-4}{J} कुछ ही कोशिकाएँ गढ़ती है जिन्हें तुरंत प्रकाशसंश्लेषण
करना ही पड़ता है।
\textbf{24.} बीज: बिना जल के निषेचन; रसद सहित, सुरक्षित, प्रसुप्त एक
ऐसा भ्रूण जो शुष्क या ऋतुबद्ध स्थानों में जम जाता है; और पंखों, फलों
तथा प्राणियों से प्रकीर्णन। \omterm{def:b2:plant-life-cycles:fern}{फ़र्न} की रणनीति
गीले, छायादार, स्थिर आवासों में जीतती है, और दूर की या नई खुली भूमि
बसाने में, जहाँ सस्ते, हलके, असंख्य प्रवर्ध्य रसद से अधिक मायने रखते
हैं।
\textbf{25.} \omterm{def:b2:plant-life-cycles:fern}{फ़र्न}: 1.9 वयस्कों के लिए प्रति ऋतु $5.1\times 10^{6}$ बीजाणु,
यानी प्रति वयस्क संतान $2.7\times 10^{6}$ बीजाणु, हर एक \qty{430}{J};
चीड़: 480 वयस्कों के लिए $9.6\times 10^{5}$ बीज, यानी प्रति
वयस्क संतान \num{2000} बीज, हर एक \qty{100}{kJ}।
\end{solution}
